\documentclass[10pt]{article}
\usepackage[top=0.75in,bottom=0.75in,left=0.75in,right=0.75in]{geometry}
\usepackage[T1]{fontenc}
\usepackage{enumitem}
\usepackage{titlesec}
\usepackage{array}
\usepackage{booktabs}

\title{AB-900: Copilot and Agent Administration Fundamentals}
\author{}
\date{}

\setcounter{secnumdepth}{2}
\setcounter{tocdepth}{2}
\renewcommand{\thesection}{\arabic{section}}
\renewcommand{\thesubsection}{\thesection.\arabic{subsection}}
\titleformat{\section}[display]{\normalfont\Huge\bfseries\centering}{\thesection}{0.5em}{\vspace*{\fill}}[\vspace*{\fill}]
\titlespacing*{\section}{0pt}{0pt}{0pt}
\titleformat{\subsection}[display]{\normalfont\LARGE\bfseries\centering}{\thesubsection}{0.5em}{\vspace*{1.25cm}}[\vspace{0.6em}\titlerule]
\titlespacing*{\subsection}{0pt}{0pt}{1.5em}

\begin{document}
\begin{titlepage}
\thispagestyle{empty}
\vspace*{\fill}
\begin{center}
{\LARGE\bfseries AB-900: Copilot and Agent Administration Fundamentals\par}
\end{center}
\vspace*{\fill}
\end{titlepage}
\tableofcontents
\newpage

\setcounter{section}{0}
\section{Microsoft 365 Security Foundations}
\setcounter{subsection}{0}
\clearpage
\subsection{Security Foundations Review}
\begin{enumerate}[leftmargin=*]

\item Which statement best describes the Zero Trust security model?
\begin{enumerate}[label=\alph*., leftmargin=1.5em]
\item Trust internal traffic by default and inspect external traffic only.
\item Never trust by default; always verify users, devices, and apps using context.
\item Focus only on firewall hardening and VPN protection.
\item Grant broad access first, then review permissions quarterly.
\end{enumerate}
\textbf{Answer:} b. Never trust by default; always verify users, devices, and apps using context.

\textbf{Explanation:} Zero Trust replaces perimeter-based trust with continuous verification of every access request, regardless of location.

\item Why is Zero Trust especially important in hybrid work environments?
\begin{enumerate}[label=\alph*., leftmargin=1.5em]
\item Because all users always work from trusted office networks.
\item Because attackers only target on-premises servers.
\item Because users, devices, and applications operate across many locations and networks.
\item Because MFA is no longer needed in hybrid work.
\end{enumerate}
\textbf{Answer:} c. Because users, devices, and applications operate across many locations and networks.

\textbf{Explanation:} Hybrid work expands the attack surface beyond the corporate perimeter, requiring adaptive and context-aware controls.

\item In Microsoft 365, Zero Trust is best understood as:
\begin{enumerate}[label=\alph*., leftmargin=1.5em]
\item A single Microsoft Defender feature.
\item A strategy spanning identity, endpoints, data, apps, infrastructure, and network.
\item A replacement for all existing governance controls.
\item A policy used only by global administrators.
\end{enumerate}
\textbf{Answer:} b. A strategy spanning identity, endpoints, data, apps, infrastructure, and network.

\textbf{Explanation:} The model is comprehensive and cross-domain, not a single product.

\item The three core principles of Zero Trust are Verify explicitly, Least privilege access, and Assume breach. (True/False)

\textbf{Answer:} True.

\textbf{Explanation:} These three principles guide access decisions, privilege boundaries, and containment strategies.

\item Which signal is \emph{not} commonly used when verifying access explicitly?
\begin{enumerate}[label=\alph*., leftmargin=1.5em]
\item Device health
n\item User behavior
\item Sign-in location
\item Current weather in the user's city
\end{enumerate}
\textbf{Answer:} d. Current weather in the user's city.

\textbf{Explanation:} Zero Trust verification uses identity, device, risk, behavior, and location signals, not unrelated environmental data.

\item A Conditional Access policy can enforce session controls such as read-only access or blocking downloads. (True/False)

\textbf{Answer:} True.

\textbf{Explanation:} Session controls help reduce data exfiltration risk in high-risk access scenarios.

\item Which scenario best reflects risk-based authentication with MFA?
\begin{enumerate}[label=\alph*., leftmargin=1.5em]
\item Prompt every sign-in equally regardless of context.
\item Never require MFA if a password is correct.
\item Trigger MFA when sign-in risk is high, such as unfamiliar location or device.
\item Allow all sign-ins from unmanaged devices.
\end{enumerate}
\textbf{Answer:} c. Trigger MFA when sign-in risk is high, such as unfamiliar location or device.

\textbf{Explanation:} Risk-based MFA balances strong security and user experience by stepping up authentication only when needed.

\item Match each concept with its best description.

\begin{center}
\begin{tabular}{>{\raggedright\arraybackslash}p{0.34\textwidth} >{\raggedright\arraybackslash}p{0.58\textwidth}}
\toprule
\textbf{Column A} & \textbf{Column B} \\
\midrule
1. RBAC & a. Aggregates logs and correlates multi-domain threats \\
2. PIM & b. Assigns predefined permissions by role \\
3. Microsoft Sentinel & c. Just-in-time elevation with approvals and audit \\
4. Defender for Endpoint & d. Detects endpoint behaviors like ransomware or lateral movement \\
\bottomrule
\end{tabular}
\end{center}

\textbf{Answer:} 1-b, 2-c, 3-a, 4-d.

\textbf{Explanation:} RBAC defines least-privilege roles, PIM limits elevated access duration, Sentinel provides SIEM/SOAR-style correlation, and Defender for Endpoint protects devices.

\item In least privilege access, users and services should receive:
\begin{enumerate}[label=\alph*., leftmargin=1.5em]
\item Maximum access to reduce admin overhead.
\item Access equal to their manager's permissions.
\item Minimum required access to perform specific tasks.
\item Permanent global administrator rights.
\end{enumerate}
\textbf{Answer:} c. Minimum required access to perform specific tasks.

\textbf{Explanation:} Least privilege shrinks attack surface and limits impact if an identity is compromised.

\item PIM primarily helps by:
\begin{enumerate}[label=\alph*., leftmargin=1.5em]
\item Removing the need for audit logs.
\item Granting permanent privileged roles.
\item Enabling temporary just-in-time privileged access.
\item Disabling Conditional Access policies.
\end{enumerate}
\textbf{Answer:} c. Enabling temporary just-in-time privileged access.

\textbf{Explanation:} PIM reduces standing privilege by elevating access only for approved, time-bound tasks.

\item The Assume breach principle emphasizes prevention only and not containment. (True/False)

\textbf{Answer:} False.

\textbf{Explanation:} Assume breach focuses on containment, segmentation, continuous monitoring, and rapid response.

\item Which Microsoft tool can isolate suspicious devices and provide endpoint forensic timelines?
\begin{enumerate}[label=\alph*., leftmargin=1.5em]
\item Microsoft Defender for Endpoint
\item Microsoft Purview DLP
\item Azure Policy
\item Microsoft Authenticator
\end{enumerate}
\textbf{Answer:} a. Microsoft Defender for Endpoint.

\textbf{Explanation:} Defender for Endpoint detects suspicious behavior and supports isolation and deep investigation.

\item Which of the following statements about Microsoft Sentinel is correct?
\begin{enumerate}[label=\alph*., leftmargin=1.5em]
\item It only collects Entra ID sign-in logs.
\item It cannot use machine learning for threat correlation.
\item It can correlate alerts across Microsoft and non-Microsoft sources and run automated playbooks.
\item It replaces all endpoint security products.
\end{enumerate}
\textbf{Answer:} c. It can correlate alerts across Microsoft and non-Microsoft sources and run automated playbooks.

\textbf{Explanation:} Sentinel centralizes detection and response by connecting events across multiple systems.

\item Which attack techniques are specifically associated with Microsoft Defender for Identity detection? \textbf{(Select all that apply)}
\begin{enumerate}[label=\alph*., leftmargin=1.5em]
\item Pass-the-Hash
\item Golden Ticket
\item Domain enumeration
\item SQL schema migration
\end{enumerate}
\textbf{Answer:} a, b, c.

\textbf{Explanation:} Defender for Identity analyzes identity-related traffic and suspicious behaviors tied to credential and directory attacks.

\item Put these Zero Trust pillars in the same set used by Microsoft 365: Identity, Endpoints, Data, Applications, Infrastructure, and \underline{\hspace{2cm}}.

\textbf{Answer:} Network.

\textbf{Explanation:} Microsoft frames Zero Trust implementation across six interdependent pillars ending with network controls.

\item Which statement best captures identity as the cornerstone pillar?
\begin{enumerate}[label=\alph*., leftmargin=1.5em]
\item Identity controls are optional if devices are compliant.
\item Identity is central because authentication and authorization decisions start there.
\item Identity applies only to external users.
\item Identity has no role in Conditional Access.
\end{enumerate}
\textbf{Answer:} b. Identity is central because authentication and authorization decisions start there.

\textbf{Explanation:} Microsoft Entra ID underpins sign-in, policy enforcement, risk detection, and access governance.

\item Which options are valid passwordless methods mentioned for Microsoft Entra ID? \textbf{(Select all that apply)}
\begin{enumerate}[label=\alph*., leftmargin=1.5em]
\item Windows Hello for Business
\item FIDO2 security keys
\item Microsoft Authenticator
\item Shared spreadsheet passphrase
\end{enumerate}
\textbf{Answer:} a, b, c.

\textbf{Explanation:} Passwordless methods reduce password risk; ad hoc shared passphrases are insecure and not approved methods.

\item Which endpoint noncompliance example is explicitly aligned with Zero Trust endpoint controls?
\begin{enumerate}[label=\alph*., leftmargin=1.5em]
\item Blocking a jailbroken phone from accessing Microsoft Teams.
\item Disabling MFA for mobile users.
\item Granting full SharePoint download rights to all unmanaged devices.
\item Ignoring minimum OS version requirements.
\end{enumerate}
\textbf{Answer:} a. Blocking a jailbroken phone from accessing Microsoft Teams.

\textbf{Explanation:} Intune compliance policies enforce device trust criteria before access is granted.

\item Which Microsoft capability addresses shadow IT and real-time app session controls?
\begin{enumerate}[label=\alph*., leftmargin=1.5em]
\item Microsoft Defender for Cloud Apps
\item Microsoft Purview Insider Risk Management
\item VPN Gateway
\item Azure Bastion
\end{enumerate}
\textbf{Answer:} a. Microsoft Defender for Cloud Apps.

\textbf{Explanation:} Defender for Cloud Apps discovers sanctioned/unsanctioned apps and can block risky behaviors in-session.

\item A document containing credit card numbers can be automatically labeled and encrypted by Microsoft Purview Information Protection. (True/False)

\textbf{Answer:} True.

\textbf{Explanation:} Purview can classify and auto-label content based on sensitive information patterns, then apply protection.

\item Which workloads are directly referenced as DLP-integrated in Microsoft 365? \textbf{(Select all that apply)}
\begin{enumerate}[label=\alph*., leftmargin=1.5em]
\item Exchange
\item SharePoint
\item OneDrive
\item Teams
\item Xbox Live
\end{enumerate}
\textbf{Answer:} a, b, c, d.

\textbf{Explanation:} DLP policy enforcement spans core Microsoft 365 collaboration and communication services.

\item Which statement about Microsoft Defender for Cloud and Azure Policy is correct?
\begin{enumerate}[label=\alph*., leftmargin=1.5em]
\item Defender for Cloud is only for Azure and can't support multicloud.
\item Azure Policy can enforce rules like tagging, encryption, and allowed resource configurations.
\item Azure Policy is a replacement for Conditional Access.
\item Defender for Cloud has no compliance relevance.
\end{enumerate}
\textbf{Answer:} b. Azure Policy can enforce rules like tagging, encryption, and allowed resource configurations.

\textbf{Explanation:} Defender for Cloud assesses posture and recommends remediation, while Azure Policy enforces governance guardrails.

\item Network controls in Zero Trust include Azure Firewall, VPN Gateway, and identity-aware monitoring for lateral movement. (True/False)

\textbf{Answer:} True.

\textbf{Explanation:} Zero Trust network strategy combines traffic control, secure remote access, and identity attack detection.

\item Order these Zero Trust response activities after suspicious activity is identified:
\begin{enumerate}[label=\alph*., leftmargin=1.5em]
\item Correlate events across domains.
\item Detect anomaly on endpoint/identity.
\item Trigger automated containment/playbooks.
\item Investigate and trace attack path.
\end{enumerate}
\textbf{Answer:} b $\rightarrow$ a $\rightarrow$ c $\rightarrow$ d.

\textbf{Explanation:} Detection comes first, followed by correlation, response automation, and deeper investigation for root cause and scope.

\item Why does Zero Trust reduce blast radius during compromise?

\textbf{Answer:} Because it limits access through least privilege, verifies context continuously, and contains threats through segmentation and rapid response.

\textbf{Explanation:} Even when an account or device is compromised, layered controls reduce lateral movement and data exposure.

\end{enumerate}

\end{document}
